\documentclass[sigconf,review,anonymous]{acmart}

\usepackage{graphicx}
\usepackage{url}
\usepackage{booktabs}
\usepackage{times}
\usepackage{amsmath}

% For double-blind review - actual IDs would be provided by CCS
\ccsSubmissionID{12345}
\ccsPaperID{1234}
\ccsISBN{978-1-4503-XXXX-X/XX/YY}

\begin{document}

\title{Ghostext: An Almost Perfect Steganography via Large Language Model and Arithmetic Coding}
\author{}
\maketitle

\begin{abstract}
We introduce Ghostext, a text steganography system that uses large language models and arithmetic coding to hide secret messages within generated text. The system allows a sender to encode a secret message along with a password and prompt into what appears to be ordinary text generated by a language model. A receiver who knows the password and prompt can recover the original message. Our approach achieves what we call almost perfect steganography because it preserves the model's original next-token probability distribution while maximizing the entropy utilization through arithmetic coding. The system combines encryption to make the message indistinguishable from random bits with arithmetic coding to efficiently use all the available entropy in the token distribution. We demonstrate that Ghostext successfully hides messages across different languages and text types while maintaining high text quality and achieving practical embedding capacity.
\end{abstract}

\keywords{steganography, large language models, arithmetic coding, information hiding, text generation}

\section{Introduction}

The ability to communicate privately is important in many situations where people need to share information without others being able to read it. Traditional encryption hides the content of messages but leaves it obvious that communication is taking place. Steganography goes a step further by hiding the very fact that a message is being sent. Text steganography has been studied for many years, but most approaches involve modifying text after it has been written, which often makes the text look unnatural or can be detected through careful analysis.

Recent advances in large language models have opened new possibilities for text steganography. These models work by predicting what word should come next based on all the words that came before, and at each step they assign probabilities to many possible next words. This probabilistic nature creates natural randomness that can be used to hide information. When the model has many reasonable next words to choose from, we can use our choice among them to encode bits of a secret message.

Ghostext takes advantage of this opportunity by combining two powerful ideas. First, we encrypt the secret message so that it looks completely random. Second, we use arithmetic coding to map this encrypted message onto the choices among the most likely next words predicted by the language model. Because we are not adding new information to the text but rather using the existing randomness in the model's predictions, the resulting text looks completely natural and follows exactly the same distribution that the model would have produced if it were just writing normally.

Our system involves three parties who we call Alice, Bob, and Censor. Alice wants to send a secret message to Bob. She provides a prompt that tells the language model what kind of text to generate, a password that only she and Bob share, and the secret message she wants to hide. The Ghostext system then produces what looks like an ordinary piece of text. When Censor examines this text, they see nothing suspicious because the text follows the same patterns that ordinary language model output follows. Censor does not have access to the prompt or the password, so they cannot tell that the text contains hidden information. Bob, who knows both the prompt and the password, can recover the original secret message.

The main contribution of our work is showing how to achieve what we call almost perfect steganography with language models. Perfect steganography would mean that the distribution of stego text is exactly the same as the distribution of ordinary text. Our approach comes close to this ideal because we carefully preserve the model's original token probability distribution while maximizing how much information we can embed. The key insight is that the combination of encryption and arithmetic coding allows us to use all the entropy that naturally exists in the language model's predictions without introducing any detectable patterns.

The rest of this paper explains how Ghostext works in detail, analyzes its security against a passive censor, describes our implementation using real language models, and presents experimental results showing that the system works well in practice.

\section{Related Work}

Steganography in the context of large language models is a relatively new but rapidly developing field. The core idea is to embed secret messages within text generated by language models, using the inherent randomness in token selection as a communication channel. Early work on text steganography predates modern language models and focused on linguistic properties such as synonym substitution, syntactic structure modification, and word ordering. These approaches typically operate on already-generated text and modify it to hide information. However, they often suffer from reduced text quality and are vulnerable to statistical analysis.

With the advent of large language models, researchers have discovered that the probabilistic nature of token selection provides a natural channel for steganography. The model produces a probability distribution over possible next tokens at each step, and this distribution contains entropy that can be exploited to encode information. One line of work focuses on using the model's sampling process directly. By carefully choosing which token to select from the top candidates based on secret message bits, information can be embedded without significantly altering the text's appearance. This approach has the advantage of preserving the model's natural language quality while providing a covert communication channel. Arithmetic coding and similar techniques have been applied to maximize the utilization of the available entropy in these distributions. By treating the secret message as a compressed bitstream and the token probability distribution as the coding model, these methods can achieve near-optimal embedding efficiency.

LLM watermarking is a closely related field that focuses on embedding detectable markers in AI-generated text rather than hiding secret messages. The goals are different---watermarking aims to identify AI-generated content, while steganography aims to hide arbitrary information---but the underlying techniques share many similarities. The most prominent approach to LLM watermarking is statistical watermarking, which modifies the token selection process to create detectable statistical patterns. A well-known example is the work by Kirchenbauer and colleagues, which introduces a watermark by splitting the vocabulary into ``green'' and ``red'' lists and biasing the selection toward green tokens. These watermarks are designed to be statistically detectable without significantly impacting the text quality. Detection typically involves analyzing the text for the expected statistical bias in token choices.

Both watermarking and steganography operate on the same principle: they manipulate the token selection process to embed information. The main difference is their purpose and threat model. Watermarking is designed to be detectable by authorized parties while being difficult to remove. Steganography is designed to be undetectable by unauthorized parties while remaining recoverable by the intended recipient. The arithmetic coding approach used in Ghostext shares conceptual similarities with some watermarking schemes, particularly in how they utilize the model's probability distribution. However, Ghostext's goal of preserving the original distribution while maximizing entropy utilization distinguishes it from most watermarking approaches, which typically introduce bias for detection purposes.

\section{System Description}

The Ghostext system takes several inputs from the user: a secret message, a password or passphrase, a cover prompt that describes what kind of text to generate, and configuration parameters that control the encoding process. The system then produces two main outputs: the stego text that contains the hidden message and information that helps verify successful decoding.

The encoding process begins with the secret message that Alice wants to send. This message is first converted to bytes using UTF-8 encoding, which ensures that any text in any language can be handled. These bytes are then encrypted using an authenticated encryption scheme. We use scrypt for key derivation, which takes the password and a randomly generated salt to produce an encryption key, and we use ChaCha20-Poly1305 as the authenticated encryption algorithm. The encryption step serves two purposes. First, it protects the confidentiality of the secret message so that only someone with the correct password can read it. Second, and more importantly from a steganography perspective, the encrypted output looks completely random, which is exactly what we need for the next step.

After encryption, we have what we call a stego packet. This packet contains not just the encrypted message but also metadata like the salt and nonces used in the encryption. The packet is formatted in a specific way that allows the receiver to know exactly how many bytes to read. This is important because the receiver needs to know when they have successfully recovered all the information without reading too much or too little.

The next part of the system is where the actual hiding happens. We treat the encrypted packet as a sequence of bits that needs to be encoded. At the same time, the language model is generating text word by word. At each step, the model looks at all the words that have been generated so far and produces a probability distribution over possible next words. This distribution tells us which words are most likely to come next.

Ghostext uses arithmetic coding to bridge the gap between the encrypted packet and the word choices. Arithmetic coding is a technique that allows us to encode a sequence of bits into a series of choices, where each choice comes from a probability distribution. The idea is to think of the encrypted packet as a number within a certain range. At each step, we look at the probability distribution that the language model provides and divide the current range into smaller ranges, one for each possible next word. We then choose the word whose range contains our target number. After making this choice, we narrow our range to just the part that corresponds to the chosen word, and we repeat the process for the next word.

The key insight is that this process does not require us to choose words that are unlikely according to the model. We can limit ourselves to only the most likely words while still making progress toward encoding our message. This means that the text we produce stays completely natural. The arithmetic coding process ensures that we use the entropy in the probability distribution efficiently. When the distribution has more uncertainty, we can encode more bits. When the distribution is more concentrated, we encode fewer bits.

The decoding process is essentially the reverse of encoding. Bob takes the stego text and runs it through the same language model with the same prompt. As the model generates each word of the text, Bob uses the same arithmetic coding process to narrow down the possible range of the encrypted packet. After seeing enough words, the range becomes small enough to determine the exact value of the packet, at which point Bob can decrypt it using the password to recover the original message.

One important design choice in Ghostext is that we only encode information at steps where the language model provides sufficient entropy. We measure the entropy of the probability distribution at each step, and if it is below a threshold, we simply choose the most likely word without encoding anything. This helps us avoid getting stuck in situations where the model is very predictable, which could cause the encoding process to stall. The system also includes safeguards to detect and retry when it encounters situations where progress becomes very slow or where certain word choices would cause tokenization problems.

The result of this process is that the stego text follows exactly the same word choice distribution that the language model would have produced if it were just writing normally. The only difference is that instead of choosing words randomly according to the distribution, we make our choices guided by the arithmetic coding process. From the perspective of someone who does not know the password, the text looks completely ordinary.

\section{Security Analysis}

The security model for Ghostext assumes a passive censor who can observe the text being communicated but does not have access to the password, the prompt, or the configuration used to generate the text. The censor can see the final text and can analyze it in any way they choose, but they cannot modify it and they do not have the shared secrets. The goal is that the censor cannot distinguish between the text that contains a hidden message and the ordinary text generated by the same language model.

Under this model, Ghostext achieves strong security through two complementary mechanisms. The first mechanism is that the language model itself produces output that is statistically similar to natural language. Because we are preserving the model's original probability distribution rather than introducing biases or patterns, any statistical properties that the censor might examine will look the same as in ordinary model output. This includes distributions of words, phrases, sentence structures, and any other linguistic features.

The second mechanism is that the encryption step makes the hidden information indistinguishable from a random bit sequence. Even if the censor somehow gained access to the intermediate values during the encoding process, the encrypted packet would provide no information about the original secret message. The use of authenticated encryption also means that any attempt to modify the stego text will be detected during decoding because the authentication tag will not verify correctly.

An important security property is that the encoding process does not introduce detectable correlations across words. Each word is chosen independently based on the probability distribution at that step and the current state of the arithmetic coder. Because the arithmetic coding process is deterministic given the encrypted packet and the sequence of probability distributions, there are no additional patterns that could be detected through correlation analysis.

The security of Ghostext can be understood through the concept of perfect steganography. Perfect steganography means that the distribution of the stego text is exactly identical to the distribution of the cover text. If this were achieved, then even a censor with unlimited computational power could not distinguish between the two. Our approach is almost perfect because we preserve the language model's original probability distribution while efficiently using the available entropy. The only difference from perfection is practical constraints like needing to encode a finite amount of information and having to handle situations where the entropy is insufficient.

From an information-theoretic perspective, the capacity of Ghostext is determined by the entropy of the language model's next-token distributions. Each token position can contribute up to the entropy of the distribution at that position, and the arithmetic coding process allows us to achieve this bound. The theoretical maximum capacity per token is the average entropy of the model's predictions, which for a well-trained language model is typically several bits per token. Our experimental results show that we achieve close to this bound in practice.

The system also includes safeguards against low-entropy regions. The language model sometimes produces very predictable output where the next word is almost certain. If we tried to encode in these regions, the encoding would either take an impractically long time or would require us to use very unlikely words that would make the text look unnatural. Ghostext detects when this happens and can automatically retry with a different salt, which changes the encrypted packet and can lead to a different path through the model's predictions.

It is important to note that Ghostext is designed against a passive censor rather than an active adversary. An active adversary who can modify the text before it reaches the receiver could potentially disrupt the hidden message. However, the authentication tag in the encryption would detect such modifications, and the receiver would know that the message has been tampered with. Protecting against active attacks that successfully recover information would require additional techniques like error-correcting codes, which we leave for future work.

\section{Technical Approach}

The technical implementation of Ghostext involves several key components that work together to achieve the steganographic goal. These components include the cryptography layer, the model backend, the candidate selection policy, the probability quantization layer, and the arithmetic coding codec.

The cryptography layer handles the encryption and decryption of the secret message. We chose ChaCha20-Poly1305 as the authenticated encryption algorithm because it provides strong security guarantees and is widely supported. The key derivation uses scrypt with parameters chosen to provide good resistance against brute-force attacks. The packet format includes a salt, a header section with metadata, and the body containing the encrypted message. The header is encrypted separately to hide information about the packet structure, which helps with the steganographic security.

The model backend is responsible for loading and running the language model. Our implementation uses llama-cpp with the Qwen3.5-2B model, which provides a good balance between quality and computational efficiency. The backend must provide two main functions: getting the next-token probability distribution for a given context and tokenizing and detokenizing text. It is critical that both the encoder and decoder use exactly the same model and tokenizer configuration, or the arithmetic coding will fail to synchronize.

The candidate selection policy determines which tokens are eligible for encoding at each step. We use a combination of top-p sampling and a maximum candidate count. The top-p parameter, set to $0.995$ in our default configuration, means we keep only the most likely tokens whose cumulative probability reaches this threshold. This helps ensure that the text stays natural while still providing sufficient entropy for encoding. The max-candidates parameter, set to $64$, ensures we do not work with too many tokens at once, which could cause tokenization instability.

An important consideration in candidate selection is retokenization stability. Sometimes a word that looks like a single choice actually consists of multiple tokens when the text is processed. We check each candidate by appending it to the current prefix, rendering the result to text, and then tokenizing again. If the token sequence changes, we exclude that candidate because the decoder might not follow the same path. This check ensures reliable decoding at the cost of excluding some tokens that would otherwise be usable.

The probability quantization layer converts floating-point probabilities from the model into integer frequencies that can be used by the arithmetic coder. We use a total frequency count of $4096$, which provides good precision while keeping the numbers manageable. The quantization process assigns integer frequencies to each candidate token such that the sum equals the total and each token gets at least one count. We use a deterministic remainder distribution to break ties, which is important for reproducibility.

The arithmetic coding codec is the core of the system. It implements a range coder that takes a message as an integer value and, given a sequence of probability distributions, produces the sequence of choices that encode that value. The encoder starts with a range covering all possible message values and narrows it down with each token choice. The decoder starts with the full range and, observing the token sequence, narrows it down until the message is uniquely determined. The implementation uses $64$-bit integers for precise arithmetic and includes renormalization to avoid underflow.

The encoding process is organized into two segments. The first segment encodes the header, which contains fixed-size metadata like the salt and nonces. The second segment encodes the body, which contains the encrypted message itself. This two-stage approach allows the decoder to first recover the header, learn the body length, and then decode the body with full knowledge of how many bytes to expect.

After the message is fully encoded, the system can optionally generate a natural tail. This tail consists of additional tokens generated without encoding anything, which helps the text end naturally rather than stopping abruptly. The tail length is bounded by a configurable maximum, and the system stops early if it detects a natural ending like a period followed by a newline.

The decoder must work with text that may contain additional trailing tokens beyond what was needed to recover the message. This is a deliberate design choice because the natural tail might be different on decoding or the text might have been modified after encoding. The decoder simply ignores any tokens after the message is fully recovered, which makes the system robust to such variations.

Error handling throughout the system is designed to fail closed rather than attempting approximate recovery. If any assumption is violated---different configurations, mismatched prompts, modified text---the system detects this and reports an error rather than producing incorrect output. This is important for security because it prevents the system from accidentally leaking partial information.

The entire pipeline is designed to be reproducible and deterministic. Given the same model, prompt, password, and configuration parameters, both encoding and decoding produce exactly the same results. This determinism is essential for the arithmetic coding to work correctly and for the system to be reliable in practice.

\section{Experiments}

We evaluated Ghostext using a real language model implementation to verify that the system works correctly and to measure its performance characteristics. Our experiments focused on three main aspects: correctness of message recovery, embedding capacity, and computational efficiency. We used the Qwen3.5-2B model with default configuration parameters including a top-p value of $0.995$, a maximum of $64$ candidates per position, and a total frequency of $4096$ for probability quantization.

\subsection{Message Recovery}

Our first set of experiments tested the ability of the system to successfully encode and recover messages across different languages and message lengths. We tested six different messages ranging from very short to moderately long, including English text, Chinese text, and even a message with numbers and special characters like coordinates. The messages varied from just two characters to over one hundred characters. For each message, we used a simple cover prompt asking the model to write a short paragraph.

The results showed perfect reliability across all test cases. Every message was successfully recovered through the round-trip process, meaning that the original secret message was identical to the decoded output. This $100\%$ success rate demonstrates that the arithmetic coding approach works correctly with real language model distributions and that the synchronization between the encoder and decoder is robust. The system handled both English and Chinese text without any issues, showing that Ghostext is effective for multiple languages.

\subsection{Embedding Capacity}

We measured the embedding capacity of the system by calculating how many bits of information were embedded per token. Across all our experiments, the bits per token ranged from approximately $2.1$ to $3.0$, with most cases falling around $2.4$ to $2.6$ bits per token. This capacity is competitive with other approaches in the literature and shows that the system makes good use of the entropy available in the language model predictions. The capacity remained stable across different message lengths, indicating that the system scales well.

The header segment of the packet, which contains metadata like salts and nonces, consistently required around $180$ to $190$ tokens to encode. This overhead is necessary for the security of the system and represents a fixed cost regardless of message length. For longer messages, this overhead becomes proportionally less significant. The body segment showed linear scaling with message length, as expected from the arithmetic coding approach.

\subsection{Computational Performance}

We also measured the computational performance of the system. The encoding speed ranged from approximately $7$ to $15$ tokens per second depending on the complexity of the generated text and the length of the message. Decoding was generally faster, typically achieving $10$ to $15$ tokens per second. These speeds are practical for interactive use, although they could be improved with more efficient hardware or model optimization. The encoding time grows roughly linearly with the number of tokens, which is expected from the sequential nature of the process.

\subsection{Retry Behavior}

One interesting observation from our experiments is the number of retry attempts. In most cases, the system succeeded on the first attempt without needing to retry. However, for some combinations of prompt and message, the system encountered low-entropy regions where progress stalled. The automatic retry mechanism successfully handled these cases by generating a new packet with fresh salt and nonces and trying again. The system is configured to allow up to ten retry attempts, and in our experiments we never needed more than four attempts even for challenging cases.

\subsection{Text Quality}

To evaluate the quality of the generated stego text, we manually inspected the outputs and verified that they appear natural and coherent. The text flows smoothly and follows the topic specified in the cover prompt. There are no obvious signs of information hiding, such as unusual word choices or disrupted sentence structure. The quality appears comparable to ordinary text generated by the same model without steganography, which is consistent with the goal of preserving the model's original distribution.

We also tested robustness to small variations in the output. Because the system includes a natural tail that can vary between encoding and decoding, we verified that adding a few extra tokens at the end of the stego text does not prevent successful decoding. The decoder simply ignores any tokens after the full packet has been recovered. This design choice makes the system more forgiving of small edits or differences in generation.

In terms of text length, we observed that the stego text length is generally proportional to the amount of information being encoded plus some overhead for the header and any natural tail. Short messages of around $10$ characters resulted in about $250$ to $300$ tokens of stego text. Longer messages of around $100$ characters resulted in about $700$ to $800$ tokens. This relationship is linear, which is expected given that each token can encode a roughly constant number of bits.

Overall, our experiments demonstrate that Ghostext successfully achieves its design goals. The system reliably hides and recovers messages across multiple languages, achieves practical embedding capacity of around $2.5$ bits per token, produces natural-looking text, and operates at reasonable computational speeds. The combination of encryption and arithmetic coding proves effective in practice, supporting the claim of almost perfect steganography that preserves the language model's original distribution while maximizing entropy utilization.

\section{Conclusion}

We have presented Ghostext, a text steganography system that achieves what we call almost perfect steganography by preserving the language model's original next-token probability distribution while efficiently encoding information through arithmetic coding. The system combines strong cryptographic security with careful use of the entropy that naturally exists in the language model predictions, allowing secret messages to be hidden within generated text that appears completely ordinary.

The key insight behind our approach is that we do not need to introduce biases or patterns into the text to hide information. Instead, we use the existing randomness in the model's token choices as a communication channel. By encrypting the secret message first, we ensure that the hidden information is indistinguishable from random bits. Then, by using arithmetic coding to map these bits onto token choices according to the model's probability distribution, we achieve efficient encoding without changing the statistical properties of the output.

Our experiments demonstrate that Ghostext works reliably in practice. The system achieved perfect message recovery across all test cases, with embedding capacity of approximately $2.5$ bits per token. The generated text maintains high quality and naturalness, making it difficult to distinguish from ordinary language model output. The computational performance is practical for interactive use, and the retry mechanism handles challenging cases gracefully.

The security of Ghostext is based on preserving the language model's original distribution, which provides strong protection against statistical analysis. Combined with cryptographic encryption and authentication, the system is secure against a passive censor who can observe the text but does not have access to the shared secrets. The system is not designed to withstand active adversaries who modify the text, but such modifications would be detected through the authentication tag.

Future work could explore several directions. One important area is improving resistance to active attacks by adding error-correcting codes that allow recovery even when the text is modified. Another direction is increasing embedding capacity through more sophisticated selection strategies that can make better use of low-entropy regions. Additionally, developing techniques for cross-model compatibility would allow the system to work across different language models and versions.

Ghostext demonstrates that high-quality text steganography is achievable with current language model technology. By carefully respecting the model's natural distribution while efficiently using available entropy, we can hide information in a way that is both practical to implement and difficult to detect. This approach opens new possibilities for secure communication in an era where large language models are increasingly used for generating text.

\begin{thebibliography}{references}
\bibitem[ACM] Kirchenbauer et al., 2023
A. Kirchenbauer, J. et al.
\newblock A Watermark for Large Language Models.
In {\em Proceedings of the 31st USENIX Security Symposium}.
USENIX Security '23,
pages 1473--1489,
year 2023.
\endbibitem

\end{thebibliography}

\end{document}
